\documentclass[11pt]{article}

\usepackage[a4paper,margin=1in]{geometry}
\usepackage{amsmath,amssymb,amsthm,mathtools}
\usepackage{enumitem}
\usepackage[hidelinks]{hyperref}

\newtheorem{theorem}{Theorem}
\newtheorem{proposition}[theorem]{Proposition}
\newtheorem{lemma}[theorem]{Lemma}
\newtheorem{corollary}[theorem]{Corollary}
\theoremstyle{definition}
\newtheorem{remark}[theorem]{Remark}
\newtheorem{problem}[theorem]{Problem}

\newcommand{\R}{\mathbb R}
\newcommand{\Om}{\Omega}
\newcommand{\Drr}{\mathcal D_{\mathrm{rr}}}
\newcommand{\Rrad}{\mathcal R_{\mathrm{rad}}}
\newcommand{\dT}{\delta_T}
\newcommand{\AF}{\mathcal A_F}
\newcommand{\Sph}{\mathbb S^{n-1}}
\newcommand{\fint}{\,\rlap{$-$}\!\!\int}
\newcommand{\norm}[1]{\left\|#1\right\|}
\newcommand{\abs}[1]{\left|#1\right|}

\title{Continuous Rearrangement and a Transport Picture for\\
       Quantitative Faber--Krahn / Saint--Venant Stability\\
       \large Toward a selection-free proof}
\author{}
\date{}

\begin{document}
\maketitle

\begin{abstract}
We develop a geometric, non-perturbative route to the sharp quantitative
Saint--Venant (equivalently Faber--Krahn) inequality that is intended to
replace the Brasco--De~Philippis--Velichkov (BDV) selection principle by a
deterministic \emph{continuous rearrangement flow}, in the spirit of the
Dolbeault--Esteban--Figalli--Frank--Loss (DEFFL) proof of sharp Sobolev
stability. The starting point is an exact splitting of the torsion deficit,
$\dT=\Drr+\Rrad$, into a \emph{non-radiality} term $\Drr$ (the Pólya--Szegő
gap of the torsion function) and a \emph{radial-profile} term $\Rrad$. The
radial-profile term is controlled by a one-dimensional spectral gap and is
elementary. The whole difficulty is concentrated in a single
\emph{quantitative Pólya--Szegő} estimate for $\Drr$, which we propose to
prove by Brock's continuous Steiner symmetrization together with a
competing-symmetries schedule, exactly as DEFFL do for Sobolev. We explain
the transport interpretation (rearrangement as a measure-preserving
deformation, dissipation as a transport cost), the monotone quantity along
the flow, and why this route never touches the regularity of $\partial\Om$
and hence bypasses selection. The remaining nucleus is stated as a precise
open problem.
\end{abstract}

\tableofcontents

\section{The transferable idea from DEFFL}

The DEFFL proof of sharp stability for the Sobolev inequality \cite{DEFFL}
runs, schematically, as follows. Let $E(f)$ be the Sobolev deficit (the
normalised gap in the inequality), whose zero set is the manifold $\mathcal
M$ of Aubin--Talenti optimisers. They construct a deformation
\[
   f \;\rightsquigarrow\; f^* \quad(\text{symmetric decreasing rearrangement}),
\]
implemented \emph{continuously} through Brock's continuous Steiner
symmetrization \cite{Brock1,Brock2}, chained over a sequence of hyperplanes
(a Bucur--Henrot schedule), and accelerated by \emph{competing symmetries}
(conformal rotations interleaved with rearrangement, after Carlen--Loss
\cite{CarlenLoss}). Two facts make this a proof of \emph{quantitative}
stability:
\begin{enumerate}[label=\textup{(\roman*)}]
\item \emph{Monotonicity / dissipation.} The Dirichlet energy is
non-increasing along the flow, and the rate at which it decreases is an
explicit non-negative \emph{dissipation} measuring how far the current
profile is from being symmetric in the active direction.
\item \emph{Local-to-global.} The flow drives any $f$ to within a
controlled distance of $\mathcal M$; one then only needs a \emph{local}
(near-optimiser) coercivity, supplied by the second variation. The deficit
controls the total dissipation, hence the total flow distance, hence the
distance of $f$ to $\mathcal M$.
\end{enumerate}
The conceptual upshot: \textbf{the continuous rearrangement flow is a
deterministic substitute for a selection/compactness step.} It transports an
arbitrary competitor to the optimiser while the deficit pays for the trip.
No penalised minimisation, no subsequence, no contradiction.

The thesis of this note is that the \emph{same} architecture should prove
sharp quantitative Saint--Venant, with the torsion function $u_\Om$ playing
the role of $f$ and the ball playing the role of the Aubin--Talenti manifold.
The role of BDV's selection principle is taken over by the rearrangement
flow.

\section{An exact rearrangement splitting of the torsion deficit}
\label{sec:splitting}

Throughout, $\Om\subset\R^n$ is bounded open, $|\Om|=\omega_n$ (the volume of
the unit ball $B=B_1$), and $u=u_\Om\in H^1_0(\Om)$ solves $-\Delta u=1$ in
$\Om$, $u=0$ on $\partial\Om$, with torsion $T(\Om)=\int_\Om u$. Write
$u^*$ for the symmetric decreasing rearrangement of $u$ (extended by $0$);
$u^*\in H^1_0(B)$ and $u^*$ is radial. Recall the convex/variational form
\[
   T(D)=\max_{v\in H^1_0(D)}J(v),\qquad J(v):=2\!\int v-\!\int|\nabla v|^2,
\]
attained at $v=u_D$, the torsion function of $D$, where $J(u_D)=\int u_D=T(D)$.

\begin{proposition}[Exact deficit splitting]\label{prop:split}
With the notation above,
\begin{equation}\label{eq:split}
   \dT(\Om)\;=\;\underbrace{\Bigl(\norm{\nabla u}_2^2-\norm{\nabla u^*}_2^2\Bigr)}_{=:\;\Drr(\Om)\;\ge\;0}
   \;+\;\underbrace{\Bigl(T(B)-J(u^*)\Bigr)}_{=:\;\Rrad(\Om)\;\ge\;0}.
\end{equation}
Here $\Drr\ge0$ is the Pólya--Szegő gap of the torsion function (vanishing
iff $u$ is, up to translation, radial decreasing), and $\Rrad\ge0$ is the
gap between the rearranged profile $u^*$ and the true ball torsion function
$u_B$ (vanishing iff $u^*=u_B$).
\end{proposition}

\begin{proof}
Since $-\Delta u=1$ and $u=0$ on $\partial\Om$, integration by parts gives
$\norm{\nabla u}_2^2=\int u=T(\Om)$. Rearrangement preserves the distribution
function, so $\int u^*=\int u$, whence $J(u^*)=2\int u-\norm{\nabla u^*}_2^2$.
Therefore
\[
   \Drr+\Rrad
   =\bigl(\norm{\nabla u}^2-\norm{\nabla u^*}^2\bigr)+\bigl(T(B)-2\!\int u+\norm{\nabla u^*}^2\bigr)
   =\norm{\nabla u}^2-2\!\int u+T(B).
\]
Using $\norm{\nabla u}^2=\int u=T(\Om)$ this equals $T(\Om)-2T(\Om)+T(B)=T(B)-T(\Om)=\dT(\Om)$.
Non-negativity of $\Drr$ is Pólya--Szegő; non-negativity of $\Rrad$ holds
because $u^*\in H^1_0(B)$ is a competitor and $T(B)=\max_{H^1_0(B)}J=J(u_B)$.
\end{proof}

\begin{remark}
The splitting \eqref{eq:split} is the exact bookkeeping behind the classical
Schwarz-symmetrization proof of Saint--Venant ($\dT\ge0$). What is new is to
read it as a sum of \emph{two independently coercive} non-negative defects
and to attack each by its own mechanism: $\Rrad$ by a one-dimensional
spectral gap (\S\ref{sec:radial}), $\Drr$ by a continuous-rearrangement
quantitative Pólya--Szegő (\S\ref{sec:dissipation}). Crucially, neither
mechanism uses any regularity of $\partial\Om$: rearrangement is defined for
arbitrary measurable competitors. This is the precise sense in which the
route is selection-free.
\end{remark}

\begin{remark}[Eigenvalue version]
The same splitting holds verbatim for the first Dirichlet eigenvalue. With
$\varphi=\varphi_\Om$ the positive first eigenfunction, $\norm\varphi_2=1$,
and $\varphi^*$ its symmetric decreasing rearrangement,
\[
   \lambda_1(\Om)-\lambda_1(B)
   =\bigl(\norm{\nabla\varphi}_2^2-\norm{\nabla\varphi^*}_2^2\bigr)
   +\bigl(R(\varphi^*)-\lambda_1(B)\bigr),
   \qquad R(\psi):=\tfrac{\norm{\nabla\psi}^2}{\norm{\psi}^2},
\]
the first term the Pólya--Szegő gap of the eigenfunction, the second the
radial Rayleigh residual. We work with torsion below because the convex
duality is cleaner; everything transfers.
\end{remark}

\section{The radial-profile term is elementary}
\label{sec:radial}

$\Rrad=J(u_B)-J(u^*)$ compares two radial functions on the fixed ball $B$.
Since $J$ is strictly concave with maximiser $u_B$, and the second variation
of $J$ at $u_B$ is $-2\int|\nabla h|^2$, restricted to radial $h\in
H^1_0(B)$, we get a clean one-dimensional coercivity.

\begin{lemma}[Radial gap]\label{lem:radial}
There is $c_n>0$ such that
\[
   \Rrad(\Om)\;\ge\;c_n\,\norm{u^*-u_B}_{H^1_0(B)}^2 .
\]
\end{lemma}

\begin{proof}
Write $u^*=u_B+h$ with $h\in H^1_0(B)$ radial and $\int(u^*-u_B)\,(-\Delta u_B-1)=0$
(the Euler--Lagrange orthogonality, automatic since $-\Delta u_B=1$). Then
$J(u^*)=J(u_B)-\int|\nabla h|^2$, so $\Rrad=\int|\nabla h|^2=\norm{\nabla h}_2^2$,
and $\norm{h}_{H^1_0}^2=\norm{\nabla h}_2^2$ by Poincaré on $B$. Hence
$\Rrad=\norm{\nabla h}_2^2\ge c_n\norm{h}_{H^1_0}^2$ with $c_n$ the Poincaré
constant of $B$.
\end{proof}

In particular, $\Rrad$ controls the deviation of the \emph{radial profile} of
$u_\Om$ from the exact ball profile. This is the easy half of the deficit and
needs nothing beyond a Poincaré inequality on the ball.

\section{The non-radiality term and the continuous flow}
\label{sec:dissipation}

The entire content of sharp stability is now in $\Drr$, the Pólya--Szegő gap.
We must show it controls the distance of $u_\Om$ from the radial cone, in a
norm strong enough to recover the Fraenkel asymmetry. This is exactly where
DEFFL's continuous machinery enters.

\subsection{Brock's continuous Steiner symmetrization}

For a direction (hyperplane) $H$, Brock's flow $\tau\mapsto u^{H}_\tau$,
$\tau\in[0,\infty]$, continuously rearranges $u$ along lines $\perp H$,
interpolating $u^H_0=u$ and $u^H_\infty=u^{*H}$ (Steiner symmetrization in
$H$). It preserves all $L^p$ norms and the distribution function, and the
Dirichlet energy is non-increasing,
\begin{equation}\label{eq:brock-mono}
   \frac{d}{d\tau}\norm{\nabla u^H_\tau}_2^2\;\le\;0,
\end{equation}
right-continuously in $\tau$ \cite{Brock1,Brock2}. Chaining a sequence of
hyperplanes $(H_k)$ with the Bucur--Henrot reparametrisation produces a
single flow $\tau\mapsto u_\tau$, $\tau\in[0,\infty]$, with $u_0=u$,
$u_\infty=u^*$, $\norm{u_\tau}_p=\norm{u}_p$, and
$\tau\mapsto\norm{\nabla u_\tau}_2^2$ non-increasing. Thus
\begin{equation}\label{eq:Drr-integral}
   \Drr(\Om)\;=\;\norm{\nabla u}_2^2-\norm{\nabla u^*}_2^2
   \;=\;\int_0^\infty\Bigl(-\tfrac{d}{d\tau}\norm{\nabla u_\tau}_2^2\Bigr)\,d\tau
   \;=\;\int_0^\infty \mathfrak d(u_\tau)\,d\tau,
\end{equation}
where $\mathfrak d(u_\tau)\ge0$ is the instantaneous \emph{dissipation}. The
flow is the deterministic ``transport to the ball'': $u$ is carried to its
radial rearrangement, and $\Drr$ is the total energy released.

\subsection{The monotone quantity and the local-to-global principle}

The monotone quantity along the flow is the Dirichlet energy
$\norm{\nabla u_\tau}_2^2$ (equivalently the running torsion functional
$J(u_\tau)=2\int u-\norm{\nabla u_\tau}^2$, non-decreasing). The DEFFL
local-to-global principle, transported to our setting, reads:

\begin{quote}
\emph{If at no time $\tau$ the profile $u_\tau$ is close to radial, then the
dissipation $\mathfrak d(u_\tau)$ stays bounded below, so $\Drr=\int\mathfrak
d\,d\tau$ is large. Contrapositively, small $\Drr$ forces $u_\tau$ to be
close to radial for a substantial range of $\tau$, in particular near
$\tau=0$, i.e.\ $u=u_\Om$ itself is close to radial.}
\end{quote}

To make this quantitative one needs a lower bound on the dissipation by a
\emph{distance to the radial cone}. Denote by
\[
   \mathcal C_{\mathrm{rad}}:=\{\,w^*(\,\cdot-z\,):z\in\R^n\,\}
\]
the cone of translated radial-decreasing rearrangements, and let
$\mathrm{dist}(u,\mathcal C_{\mathrm{rad}})$ be measured in a norm to be
specified. The nucleus of the whole program is:

\begin{problem}[Quantitative Pólya--Szegő for the torsion function]\label{prob:QPS}
Find a norm $\norm{\cdot}_X$ controlling the Fraenkel asymmetry of the
super-level sets and a constant $c=c(n)>0$ such that
\begin{equation}\label{eq:QPS}
   \Drr(\Om)\;=\;\norm{\nabla u}_2^2-\norm{\nabla u^*}_2^2
   \;\ge\;c\,\inf_{z\in\R^n}\,\norm{\,u-u^*(\,\cdot-z\,)\,}_X^2 .
\end{equation}
\end{problem}

Two features of \eqref{eq:QPS} must be respected, and they are exactly the
features DEFFL engineer for Sobolev.

\paragraph{(a) Degeneracy of one-shot Pólya--Szegő.} The bare gap
$\norm{\nabla u}^2-\norm{\nabla u^*}^2$ is \emph{not} coercive in general:
Cianchi--Fusco exhibit families where the gap is much smaller than the
asymmetry, because a single rearrangement can be ``almost free'' while still
moving mass. This is the rearrangement analogue of the obstruction we met
elsewhere in this project. The cure, after DEFFL, is to not use one
rearrangement but the \emph{continuous flow}: integrate the dissipation
\eqref{eq:Drr-integral} and bound it below using the spectral structure of
the flow at \emph{every} time, not just the endpoints.

\paragraph{(b) Competing symmetries to prevent stalling.} A single Steiner
direction can leave the profile invariant (if $u$ is already symmetric in
that direction) without being radial. DEFFL interleave conformal rotations
so that the schedule cannot stall away from the optimiser. The Faber--Krahn
analogue of the conformal group is absent (the torsion problem is not
conformally invariant), so a \emph{different} anti-stalling device is needed:
the natural candidate is to randomise/cycle the Steiner directions densely
(a.e.\ schedule converges, by \cite{VanSchaftingen}) and to use the
\emph{rotational averaging} of the dissipation. Establishing a uniform
lower bound on the time-averaged dissipation in terms of non-radiality is the
technical heart.

\subsection{The instantaneous dissipation, explicitly}

For a single Steiner direction $H=e^\perp$, slicing $\R^n=\R_e\times H$ and
writing $u(s,y)$ ($s\in\R$, $y\in H$), Brock's dissipation has the
one-dimensional-slice form
\begin{equation}\label{eq:diss-slice}
   \mathfrak d_H(u)\;=\;\int_H \mathfrak d_{1}\bigl(u(\cdot,y)\bigr)\,dy,
\end{equation}
where $\mathfrak d_1(g)\ge0$ measures, for the one-dimensional profile $g$,
the rate of Dirichlet-energy loss under the 1D continuous symmetrization. For
a smooth slice with super-level intervals $\{g>t\}=\bigcup_j(a_j(t),b_j(t))$,
$\mathfrak d_1$ is a positive quadratic form in the velocities of the
interval endpoints, vanishing exactly when the intervals are concentric. The
key elementary fact is that $\mathfrak d_1(g)$ controls the
\emph{1D asymmetry} of $g$: how far the super-level sets of $g$ are from being
centred intervals. Integrating over slices $y\in H$ and over directions
$H\in\Sph$, the dissipation controls the full (non-radial) asymmetry. Turning
this chain of ``controls'' into a quantitative inequality with a uniform
constant is Problem~\ref{prob:QPS}.

\section{Transport interpretation}
\label{sec:transport}

The rearrangement $u\mapsto u^*$ is a measure-preserving transport at the
level of super-level sets: for each height $t$, the set $\{u>t\}$ is
transported to the ball $\{u^*>t\}=B_{r(t)}$ of equal volume. The continuous
flow $u_\tau$ realises this transport through a one-parameter family of
measure-preserving maps $T_\tau$, with $T_0=\mathrm{id}$ and $(T_\infty)_\#$
sending each level set to its centred ball. In this picture:
\begin{itemize}[leftmargin=2em]
\item the \emph{monotone quantity} is the Dirichlet energy $\int|\nabla
u_\tau|^2$, decreasing along transport;
\item the \emph{dissipation} $\mathfrak d(u_\tau)$ is the instantaneous
transport cost (the squared transport velocity weighted by the gradient);
\item the \emph{total cost} $\int_0^\infty\mathfrak d\,d\tau=\Drr$ is the
deficit's payment for moving $u$ to the ball.
\end{itemize}
This is the precise mirror of the isoperimetric ``transport to the ball''
heuristic, now at the level of the torsion function rather than the bare
domain. The advantage over a raw domain transport is that the gradient weight
in the cost makes spectrally inefficient features (thin tentacles, spikes,
high-frequency boundary oscillation) \emph{cheap to remove and yet large in
energy}, which is exactly the favourable imbalance we verified directly via
the Dirichlet-energy scaling of thin features in the companion
energy-identity analysis.

\begin{remark}[Relation to the pressure flow]
The companion note develops a \emph{domain} flow with normal velocity
$V_n=(\partial_\nu u)^2-\fint(\partial_\nu u)^2$ (boundary-pressure
equalisation), whose equilibrium is the ball by Serrin. That flow is
intrinsically perturbative: it requires a smooth moving boundary, and is the
infinitesimal/Hadamard picture. The rearrangement flow of this note is its
\emph{non-perturbative} counterpart: it acts on the function, not the
boundary, is defined for arbitrary rough $\Om$, and integrates a global
dissipation. The two are linked by the coarea formula --- the pressure
$(\partial_\nu u)^2$ is the boundary trace of $|\nabla u|^2$, whose
super-level redistribution is what the rearrangement flow performs --- but
only the rearrangement flow survives at the rough, global scale where
selection was previously needed.
\end{remark}

\section{Recovering the Fraenkel asymmetry}
\label{sec:asymmetry}

Granting Problem~\ref{prob:QPS} in a norm $X$ controlling super-level
asymmetry, the deficit controls the domain asymmetry as follows. By the
triangle inequality and Proposition~\ref{prop:split},
\[
   \inf_z\norm{u-u_B(\,\cdot-z\,)}_X
   \;\le\;\inf_z\norm{u-u^*(\,\cdot-z\,)}_X+\norm{u^*-u_B}_X
   \;\lesssim\;\Drr^{1/2}+\Rrad^{1/2}\;\lesssim\;\dT^{1/2}.
\]
The super-level sets of $u$ are then close to those of $u_B(\cdot-z)$,
uniformly in height. Two points remain, both standard in flavour but
requiring care at the rough scale:
\begin{enumerate}[label=\textup{(\roman*)}]
\item \emph{From level sets to the domain.} The Fraenkel asymmetry is the
$t\to0$ level $\{u>0\}=\Om$. Control of all positive levels plus the Hopf
non-degeneracy $u\sim\mathrm{dist}(\cdot,\partial\Om)$ near the boundary
upgrades level-set control to $\AF(\Om)\lesssim\dT^{1/2}$. Unlike the
single-good-level approach, the rearrangement controls \emph{all} levels
simultaneously, which should make the $t\to0$ passage more robust.
\item \emph{Sharp exponent.} The target is $\AF(\Om)^2\le C\,\dT(\Om)$,
i.e.\ exponent $2$. The splitting and Lemma~\ref{lem:radial} are already at
the sharp quadratic scale; the exponent in the final estimate is inherited
entirely from the exponent in \eqref{eq:QPS}. Thus \emph{sharpness of the
quantitative Pólya--Szegő is equivalent to sharpness of Faber--Krahn} along
this route --- a clean reduction.
\end{enumerate}

\section{Why this bypasses selection}

\begin{enumerate}[label=\textup{(\arabic*)},leftmargin=2em]
\item \emph{No regularity of $\partial\Om$ is ever used.} Rearrangement and
Brock's flow are defined for arbitrary measurable $u$; the dissipation is a
function-space quantity. The selection step existed only to manufacture a
smooth near-equilibrium boundary on which BNST / second variation could run;
here there is no boundary analysis at all until the very end
(\S\ref{sec:asymmetry}(i)), and even there only Hopf non-degeneracy is
needed.
\item \emph{No compactness or contradiction.} The flow is explicit and
deterministic, and $\dT=\Drr+\Rrad$ is an identity. Stability is a
\emph{computation} of how the deficit is distributed between non-radiality
and profile, exactly as DEFFL turn the Sobolev deficit into an explicit flow
budget.
\item \emph{Global and non-perturbative.} The estimate \eqref{eq:QPS}, if
true with a uniform constant, holds for all $\Om$, not just near the ball.
The perturbative second-variation analysis (our $(G)$, BDV Thm 3.3) is
recovered as the \emph{linearisation} of \eqref{eq:QPS} at $u^*=u_B$, but is
no longer the load-bearing input.
\end{enumerate}

\section{The nucleus, and an honest assessment}

The program reduces sharp quantitative Saint--Venant, with no selection, to:

\begin{quote}
\textbf{Problem~\ref{prob:QPS}, sharp form.} Prove
$\norm{\nabla u}_2^2-\norm{\nabla u^*}_2^2\ge c(n)\,\inf_z\,\big\|\,
\{u>\cdot\}\,\triangle\,\{u^*(\cdot-z)>\cdot\}\,\big\|^2$ in a level-set
asymmetry norm at the quadratic scale, by integrating Brock's dissipation
\eqref{eq:Drr-integral} along a non-stalling Steiner schedule.
\end{quote}

What is solid:
\begin{itemize}[leftmargin=2em]
\item the exact splitting $\dT=\Drr+\Rrad$ (Proposition~\ref{prop:split}),
both terms non-negative, no regularity used;
\item the radial half $\Rrad\gtrsim\norm{u^*-u_B}^2_{H^1}$
(Lemma~\ref{lem:radial}), elementary;
\item the reduction of the \emph{exponent} to that of the quantitative
Pólya--Szegő;
\item the structural match with DEFFL: monotone Dirichlet energy, integrated
dissipation, local-to-global via the flow.
\end{itemize}

What is open, and genuinely hard:
\begin{itemize}[leftmargin=2em]
\item Problem~\ref{prob:QPS} itself. One-shot Pólya--Szegő is known to
degenerate (Cianchi--Fusco); the continuous flow is the proposed cure but the
dissipation lower bound \eqref{eq:diss-slice} must be made quantitative and
uniform.
\item The \emph{anti-stalling} device. The torsion problem lacks the
conformal symmetry that DEFFL exploit via competing symmetries; a substitute
(dense/random Steiner schedule with a time-averaged dissipation bound) must
be supplied. This is the single place where the Sobolev proof does not
transfer verbatim, and where new ideas are required.
\item The $t\to0$ passage from level-set control to domain asymmetry at the
rough scale (\S\ref{sec:asymmetry}(i)).
\end{itemize}

\paragraph{Verdict.} The rearrangement/transport route gives a genuinely
non-perturbative architecture for quantitative Faber--Krahn in which the
selection principle is replaced by a deterministic flow and the entire
difficulty is isolated in one quantitative Pólya--Szegő inequality. It is the
most plausible candidate among the routes considered in this project for a
proof with real geometric content that does not pass through BDV selection.
Its viability hinges on Problem~\ref{prob:QPS} and on finding the right
anti-stalling substitute for competing symmetries; both are concrete and, in
the Sobolev case, were solved exactly by DEFFL --- which is the strongest
available evidence that the torsion case is within reach.

\begin{thebibliography}{9}

\bibitem{DEFFL} J.~Dolbeault, M.~J.~Esteban, A.~Figalli, R.~L.~Frank, and
M.~Loss, \emph{Sharp stability for Sobolev and log-Sobolev inequalities,
with optimal dimensional dependence}, arXiv:2209.08651v5 (2025).

\bibitem{Brock1} F.~Brock, \emph{Continuous Steiner-symmetrization},
Math.~Nachr.\ \textbf{172} (1995), 25--48.

\bibitem{Brock2} F.~Brock, \emph{Continuous rearrangement and symmetry of
solutions of elliptic problems}, Proc.~Indian Acad.~Sci.\ Math.~Sci.\
\textbf{110} (2000), 157--204.

\bibitem{CarlenLoss} E.~A.~Carlen and M.~Loss, \emph{Extremals of
functionals with competing symmetries}, J.~Funct.~Anal.\ \textbf{88} (1990),
437--456.

\bibitem{VanSchaftingen} J.~Van Schaftingen, \emph{Universal approximation of
symmetrizations by polarizations}, Proc.~Amer.~Math.~Soc.\ \textbf{134}
(2006), 177--186.

\bibitem{BDV} L.~Brasco, G.~De~Philippis, and B.~Velichkov, \emph{Faber--Krahn
inequalities in sharp quantitative form}, Duke Math.~J.\ \textbf{164}
(2015), 1777--1831.

\bibitem{CianchiFusco} A.~Cianchi and N.~Fusco, \emph{Functions of bounded
variation and rearrangements}, Arch.~Ration.~Mech.~Anal.\ \textbf{165}
(2002), 1--40.

\end{thebibliography}

\end{document}
